\subsection{The Rabin-Karp Algorithm}

\subsubsection{Introduction \& Application}
The Rabin-Karp algorithm is a string-matching algorithm that utilizes a hash function to find occurrences of a pattern string $P$ within a text string $T$. Instead of comparing the pattern character by character against every possible window in the text, it computes a hash value for the pattern and a rolling hash value for each window of the text. If the hash values match, it performs a secondary character-by-character comparison to ensure it is not a hash collision. This rolling hash technique allows the algorithm to update the hash value of the next window in $\mathcal{O}(1)$ time by dropping the first character of the previous window and appending the new character.

\textbf{Application:} Because Rabin-Karp calculates hash values, it is exceptionally efficient when searching for \textit{multiple patterns} simultaneously. In real-world applications, it is widely used in plagiarism detection systems (where multiple sentences or phrases need to be matched against a massive database of documents) and in computational biology for DNA sequence matching.

\subsubsection{Pseudocode}
The following pseudocode demonstrates the 1D Rabin-Karp matcher. In our C++ implementation for the crossword grid, we utilize $d = 256$ (the number of characters in the extended ASCII set) and a large prime $q = 1000000007$ to minimize hash collisions. We also explicitly handle negative modulo results when updating the hash (as seen in lines 16-18).

\begin{algorithm}[H]
\caption{Rabin-Karp-Matcher(T, P, d, q)}
\begin{algorithmic}[1]
\STATE $n \leftarrow \mathrm{length}[T]$
\STATE $m \leftarrow \mathrm{length}[P]$
\STATE $h \leftarrow d^{m-1} \bmod q$
\STATE $p \leftarrow 0$ \COMMENT{Hash value for pattern}
\STATE $t_0 \leftarrow 0$ \COMMENT{Hash value for text window}
\FOR{$i \leftarrow 1$ \textbf{to} $m$}
    \STATE $p \leftarrow (d \cdot p + P[i]) \bmod q$
    \STATE $t_0 \leftarrow (d \cdot t_0 + T[i]) \bmod q$
\ENDFOR
\FOR{$s \leftarrow 0$ \textbf{to} $n - m$}
    \IF{$p == t_s$}
        \IF{$P[1..m] == T[s+1 .. s+m]$}
            \STATE \textbf{print} ``Pattern occurs with shift '' $s$
        \ENDIF
    \ENDIF
    \IF{$s < n - m$}
        \STATE $t_{s+1} \leftarrow (d(t_s - T[s+1] \cdot h) + T[s+m+1]) \bmod q$
        \IF{$t_{s+1} < 0$}
            \STATE $t_{s+1} \leftarrow t_{s+1} + q$
        \ENDIF
    \ENDIF
\ENDFOR
\end{algorithmic}
\end{algorithm}
\subsubsection{Step-by-step Trace}
We trace the algorithm using the example defined in our Problem Statement: $T = \text{``ababcab''}$ and $P = \text{``abc''}$. 
To simplify the mathematical trace, let us map the characters to integers $(a=1, b=2, c=3)$, use a small alphabet size $d=10$, and a prime number $q=13$.
\begin{itemize}
    \item \textbf{Pre-computation:} 
    \begin{itemize}
        \item $h = d^{m-1} \bmod q = 10^2 \bmod 13 = 100 \bmod 13 = 9$.
        \item $P = \text{``abc''} \rightarrow 1, 2, 3$. Hash $p = (1 \cdot 100 + 2 \cdot 10 + 3) \bmod 13 = 123 \bmod 13 = 6$.
    \end{itemize}
    \item \textbf{Window Shifts:}
    \begin{itemize}
        \item \textbf{Shift $s=0$:} $T[1..3] = \text{``aba''} \rightarrow 1, 2, 1$. \\
        $t_0 = 121 \bmod 13 = 4$. \\
        $t_0 \neq p$ ($4 \neq 6$). No match.
        
        \item \textbf{Shift $s=1$:} $T[2..4] = \text{``bab''} \rightarrow 2, 1, 2$. \\
        $t_1 = (10 \cdot (t_0 - T[1] \cdot h) + T[4]) \bmod q = (10 \cdot (4 - 1 \cdot 9) + 2) \bmod 13 = -48 \bmod 13 \equiv 4$. \\
        $t_1 \neq p$ ($4 \neq 6$). No match.
        
        \item \textbf{Shift $s=2$:} $T[3..5] = \text{``abc''} \rightarrow 1, 2, 3$. \\
        $t_2 = (10 \cdot (t_1 - T[2] \cdot h) + T[5]) \bmod 13 = (10 \cdot (4 - 2 \cdot 9) + 3) \bmod 13 = -137 \bmod 13 \equiv 6$. \\
        $t_2 == p$ ($6 == 6$). We perform a character comparison: $\text{``abc''} == \text{``abc''}$. \textbf{Match found at shift 2!}
        
        \item \textbf{Shift $s=3$:} $T[4..6] = \text{``bca''} \rightarrow 2, 3, 1$. \\
        $t_3 = (10 \cdot (t_2 - T[3] \cdot h) + T[6]) \bmod 13 = (10 \cdot (6 - 1 \cdot 9) + 1) \bmod 13 = -29 \bmod 13 \equiv 10$. \\
        $t_3 \neq p$ ($10 \neq 6$). No match.
        
        \item \textbf{Shift $s=4$:} $T[5..7] = \text{``cab''} \rightarrow 3, 1, 2$. \\
        $t_4 = (10 \cdot (t_3 - T[4] \cdot h) + T[7]) \bmod 13 = (10 \cdot (10 - 2 \cdot 9) + 2) \bmod 13 = -78 \bmod 13 \equiv 0$. \\
        $t_4 \neq p$ ($0 \neq 6$). No match.
    \end{itemize}
\end{itemize}

\subsubsection{Complexity}
\begin{itemize}
    \item \textbf{Time Complexity:} 
    \begin{itemize}
        \item \textit{Pre-processing:} $\mathcal{O}(m)$ time to compute the hash of the pattern and the first window.
        \item \textit{Matching (Best/Average Case):} $\mathcal{O}(n - m + 1)$. The rolling hash computes the next window's hash in $\mathcal{O}(1)$ time. If collisions are rare, the total time is $\mathcal{O}(n + m)$. For our crossword implementation searching a $R \times C$ grid with $K$ keywords, the average time is $\mathcal{O}(R \cdot C \cdot K)$.
        \item \textit{Worst Case:} $\mathcal{O}(n \cdot m)$. This occurs when the hash function is poor, resulting in a spurious hit (hash collision) at every single shift, forcing the algorithm to perform $m$ character comparisons each time.
    \end{itemize}
    \item \textbf{Space Complexity:} $\mathcal{O}(1)$ auxiliary space, as it only requires a few variables to store the hash values and prime modulo $q$.
\end{itemize}

\subsubsection{Discussion}
\textbf{When is it good?} The Rabin-Karp algorithm shines when we need to search for multiple patterns simultaneously (by pre-calculating the hashes of all patterns). It is highly efficient in practice when a large prime $q$ is chosen, minimizing the chance of hash collisions.

\textbf{When is it bad?} The algorithm degrades gracefully to the Naive Brute-Force approach in the worst-case scenario where hash collisions are frequent. Furthermore, the arithmetic operations involving modulo $q$ can be slower than the simple pointer increments used in algorithms like KMP or Boyer-Moore.